\documentclass[a4paper,12pt]{article}

\usepackage[bahasa]{babel}
\usepackage{geometry}
\usepackage{graphicx}
\usepackage{hyperref}
\usepackage{enumitem}

\geometry{margin=2.5cm}

\title{\textbf{Jakarta: Ibu Kota Indonesia}}
\author{}
\date{}

\begin{document}

\maketitle
\thispagestyle{empty}

\section{Pendahuluan}

Jakarta adalah ibu kota negara Indonesia dan merupakan salah satu kota metropolitan terbesar di dunia. Terletak di pesisir barat laut Pulau Jawa, Jakarta memiliki luas sekitar 661 km\textsuperscript{2} dan dihuni oleh lebih dari 10,5 juta jiwa.

\section{Sejarah Singkat}

Jakarta memiliki sejarah panjang yang dimulai dari abad ke-4:
\begin{itemize}
    \item \textbf{Sunda Kelapa} --- Pelabuhan utama Kerajaan Sunda
    \item \textbf{Jayakarta} --- Nama yang diberikan oleh Fatahillah pada tahun 1527
    \item \textbf{Betavia} --- Nama selama masa penjajahan Belanda
    \item \textbf{Jakarta} --- Nama resmi setelah kemerdekaan Indonesia
\end{itemize}

\section{Geografi dan Demografi}

Jakarta berbatasan dengan:
\begin{description}
    \item[Utara] Laut Jawa
    \item[Selatan] Kota Depok dan Kabupaten Bogor
    \item[Barat] Kota Tangerang
    \item[Timur] Kota Bekasi
\end{description}

Provinsi DKI Jakarta terdiri dari lima kota administrasi dan satu kabupaten administrasi, yaitu:
\begin{enumerate}
    \item Jakarta Pusat
    \item Jakarta Utara
    \item Jakarta Barat
    \item Jakarta Selatan
    \item Jakarta Timur
    \item Kepulauan Seribu
\end{enumerate}

\section{Ekonomi dan Transportasi}

Sebagai pusat bisnis dan keuangan Indonesia, Jakarta menyumbang produk domestik bruto (PDB) nasional yang signifikan. Beberapa moda transportasi utama di Jakarta meliputi:

\begin{itemize}
    \item TransJakarta (Bus Rapid Transit)
    \item MRT Jakarta (Moda Raya Terpadu)
    \item LRT Jakarta (Lintas Raya Terpadu)
    \item KRL Commuter Line
\end{itemize}

\section{Masalah dan Tantangan}

Beberapa tantangan utama yang dihadapi Jakarta antara lain:
\begin{itemize}
    \item Kemacetan lalu lintas
    \item Banjir musiman
    \item Penurunan permukaan tanah
    \item Polusi udara
\end{itemize}

Pada tahun 2024, pemerintah Indonesia resmi memindahkan ibu kota negara ke Nusantara di Kalimantan Timur, namun Jakarta tetap menjadi pusat ekonomi dan bisnis utama Indonesia.

\begin{thebibliography}{1}
    \bibitem{bps} Badan Pusat Statistik. \textit{Jakarta dalam Angka 2024}.
    \bibitem{history} Ricklefs, M. C. \textit{Sejarah Indonesia Modern}. Yogyakarta: Gadjah Mada University Press, 2005.
\end{thebibliography}

\end{document}
